\documentclass[12pt,a4paper]{article}
\usepackage[UTF8]{ctex}
\usepackage{amsmath,amssymb}
\usepackage{geometry}
\usepackage{hyperref}
\usepackage{listings}
\usepackage{xcolor}
\usepackage{booktabs}
\usepackage{graphicx}

\geometry{left=2.5cm,right=2.5cm,top=2.5cm,bottom=2.5cm}

% 代码样式
\lstdefinestyle{pythonstyle}{
    language=Python,
    basicstyle=\ttfamily\small,
    keywordstyle=\color{blue}\bfseries,
    stringstyle=\color{red},
    commentstyle=\color{green!60!black}\itshape,
    numbers=left,
    numberstyle=\tiny\color{gray},
    frame=single,
    breaklines=true,
    showstringspaces=false,
    tabsize=4,
    backgroundcolor=\color{gray!5},
    rulecolor=\color{gray!30}
}
\lstset{style=pythonstyle}

\title{\textbf{JAX 框架入门介绍}}
\author{面向机器学习研究者与工程师}
\date{\today}

\begin{document}

\maketitle

\begin{abstract}
JAX 是由 Google Research 开发的高性能数值计算框架，它将 NumPy 的简洁接口、自动微分的能力和硬件加速的性能融为一体。本文简要介绍 JAX 的核心设计理念、四大函数变换原语及其典型应用场景，帮助读者快速上手这一在机器学习研究领域日益流行的工具。
\end{abstract}

\tableofcontents
\newpage

% =============================================
\section{什么是 JAX？}
% =============================================

JAX 的全称可以理解为 \textbf{NumPy on Accelerators}。它的核心定位可以用一句话概括：

\begin{quote}
\emph{JAX = NumPy + 自动微分 + GPU/TPU 加速编译}
\end{quote}

与 PyTorch、TensorFlow 等主流深度学习框架相比，JAX 的设计哲学有显著不同：它不是一个"大而全"的训练框架，而是一组\textbf{可组合的函数变换工具}。用户以纯函数的形式定义计算逻辑，JAX 则提供一套高阶函数来对这些计算进行变换——求导、向量化、编译加速、并行化——就像搭积木一样自由组合。

\subsection{核心优势}

\begin{itemize}
    \item \textbf{NumPy 兼容}：\texttt{jax.numpy} 提供了与 NumPy 几乎一致的 API，迁移成本极低。
    \item \textbf{自动微分}：原生支持前向模式（JVP）和反向模式（VJP），可轻松计算高阶导数。
    \item \textbf{即时编译（JIT）}：通过 XLA 编译器将 Python 代码编译为高效的 GPU/TPU 内核。
    \item \textbf{向量化}：\texttt{vmap} 自动将单样本函数提升为批处理函数，无需手动编写循环。
    \item \textbf{纯函数式设计}：无副作用、不可变数据，使程序变换在语义上安全且可预测。
\end{itemize}

% =============================================
\section{安装与快速上手}
% =============================================

\subsection{安装}

JAX 的安装非常简单，可通过 pip 完成。以 CPU 版本为例：

\begin{lstlisting}
pip install jax jaxlib
\end{lstlisting}

若需 GPU 支持，请参考官方文档安装对应的 CUDA 版本。

\subsection{第一个例子}

下面是一个最简单的 JAX 程序——计算函数 $f(x) = x^2$ 在 $x=3.0$ 处的导数：

\begin{lstlisting}
import jax
import jax.numpy as jnp

# 定义函数
def f(x):
    return x ** 2

# 使用 jax.grad 自动求导
df = jax.grad(f)

# 计算 f'(3.0) = 2 * 3.0 = 6.0
print(df(3.0))  # 输出: 6.0
\end{lstlisting}

\texttt{jax.grad} 接受一个标量值函数，返回其梯度函数。这一设计体现了 JAX 的核心理念：\textbf{函数变换}——输入一个函数，输出另一个函数。

% =============================================
\section{四大核心变换}
% =============================================

JAX 的能力建立在四个核心变换之上，它们可以\textbf{任意组合}使用。

\subsection{grad --- 自动微分}

\texttt{grad} 用于计算标量函数关于其参数的梯度，底层基于反向模式自动微分（即反向传播）。

\begin{lstlisting}
def loss(w, b, x, y):
    pred = w * x + b
    return (pred - y) ** 2

# 分别对 w 和 b 求偏导
grad_w = jax.grad(loss, argnums=0)  # 第0个参数
grad_b = jax.grad(loss, argnums=1)  # 第1个参数

print(grad_w(2.0, 1.0, 3.0, 4.0))  # d(loss)/dw
print(grad_b(2.0, 1.0, 3.0, 4.0))  # d(loss)/db
\end{lstlisting}

\textbf{高阶导数}同样轻而易举——只需对 \texttt{grad} 反复应用：

\begin{lstlisting}
f = lambda x: x ** 3
df = jax.grad(f)           # 一阶导: 3x^2
d2f = jax.grad(df)         # 二阶导: 6x
d3f = jax.grad(d2f)        # 三阶导: 6

print(d2f(2.0))  # 输出: 12.0
\end{lstlisting}

\subsection{jit --- 即时编译加速}

\texttt{jit}（Just-In-Time compilation）通过 XLA 编译器将函数编译为优化的机器码，消除 Python 解释器的开销：

\begin{lstlisting}
@jax.jit
def heavy_computation(x):
    for _ in range(1000):
        x = x * 2
    return x

# 首次调用时编译，后续调用直接执行编译后的代码
result = heavy_computation(jnp.array(1.0))
\end{lstlisting}

在实际的深度学习训练中，\texttt{jit} 可以将训练步骤的执行速度提升数倍至数十倍。

\subsection{vmap --- 自动向量化}

\texttt{vmap} 将一个处理单个样本的函数自动提升为处理一批样本的函数，无需手动编写批处理逻辑：

\begin{lstlisting}
# 单个样本的预测
def predict(params, x):
    w, b = params
    return jnp.dot(w, x) + b

# 自动向量化为批处理
batched_predict = jax.vmap(predict, in_axes=(None, 0))

# 对一批输入进行预测
batch_x = jnp.ones((32, 784))  # 32 个样本
results = batched_predict(params, batch_x)  # 形状: (32,)
\end{lstlisting}

\texttt{in\_axes} 参数指定每个输入沿哪个轴进行映射：\texttt{None} 表示不映射（广播），\texttt{0} 表示沿第 0 轴映射。

\subsection{pmap --- 并行映射}

\texttt{pmap} 与 \texttt{vmap} 类似，但将计算分配到多个 GPU/TPU 设备上并行执行：

\begin{lstlisting}
# 假设有 8 个 GPU
parallel_train = jax.pmap(train_step)

# 数据自动分配到 8 个设备上并行训练
results = parallel_train(batched_data)
\end{lstlisting}

\subsection{变换的组合}

JAX 最强大的特性在于这些变换可以\textbf{自由嵌套}。例如：

\begin{lstlisting}
# 编译加速 + 向量化 + 求导
fast_batch_grad = jax.jit(jax.vmap(jax.grad(loss)))
\end{lstlisting}

这一行代码的含义是：对损失函数求梯度，然后向量化到一批数据上，最后编译加速。整个过程无需手动处理任何底层细节。

% =============================================
\section{JAX 的纯函数式约束}
% =============================================

JAX 的变换（\texttt{jit}, \texttt{grad}, \texttt{vmap} 等）要求被变换的函数是\textbf{纯函数}（pure function），即：

\begin{itemize}
    \item 输出仅依赖于输入参数，不依赖外部全局状态。
    \item 不产生副作用（如修改全局变量、打印输出等）。
\end{itemize}

\subsection{常见的"坑"与解决方案}

\subsubsection{不可变数组}

JAX 数组（\texttt{jnp.ndarray}）是\textbf{不可变的}：

\begin{lstlisting}
x = jnp.array([1, 2, 3])
# x[0] = 10  # 错误！JAX 数组不支持原地修改

# 正确做法：使用 .at 语法
x = x.at[0].set(10)
print(x)  # [10, 2, 3]
\end{lstlisting}

\subsubsection{显式随机数管理}

JAX 不使用全局随机状态，而是要求显式传递 PRNG key：

\begin{lstlisting}
key = jax.random.PRNGKey(42)

# 每次使用需要拆分 key
key, subkey = jax.random.split(key)
sample = jax.random.normal(subkey, shape=(3,))

key, subkey = jax.random.split(key)
another_sample = jax.random.normal(subkey, shape=(3,))
\end{lstlisting}

这种设计虽然略显繁琐，但确保了随机数的\textbf{可复现性}和\textbf{并行安全性}。

% =============================================
\section{JAX 生态与上层库}
% =============================================

JAX 本身是一个底层计算框架，社区在此基础上构建了丰富的上层库：

\begin{table}[htbp]
\centering
\caption{JAX 生态系统中的主要上层库}
\label{tab:ecosystem}
\begin{tabular}{@{}lll@{}}
\toprule
\textbf{库名} & \textbf{用途} & \textbf{类比} \\
\midrule
Flax        & 神经网络模块化       & PyTorch nn.Module \\
Haiku       & 神经网络模块化（DeepMind） & PyTorch nn.Module \\
Optax       & 优化器与梯度变换     & torch.optim       \\
Optax       & 优化器与梯度变换     & torch.optim       \\
Equinox     & 基于 PyTree 的模块化 & nn.Module + JAX   \\
Diffrax     & 可微分 ODE/SDE 求解器 & torchdiffech      \\
Orbax       & 检查点与模型序列化   & torch.save/load   \\
\bottomrule
\end{tabular}
\end{table}

% =============================================
\section{JAX 与 PyTorch 的对比}
% =============================================

\begin{table}[htbp]
\centering
\caption{JAX 与 PyTorch 的核心差异}
\label{tab:comparison}
\begin{tabular}{@{}lll@{}}
\toprule
\textbf{维度}       & \textbf{JAX}              & \textbf{PyTorch}           \\
\midrule
编程范式       & 纯函数式                 & 面向对象 + 命令式           \\
数组可变性     & 不可变                   & 可变（原地操作）             \\
编译加速       & 内置 JIT（XLA）           & torch.compile（较新）       \\
向量化         & vmap（自动）              & 手动编写批处理逻辑            \\
自动微分       & 前向 + 反向 + 高阶       & 主要为反向模式               \\
随机数管理     & 显式 PRNG key            & 隐式全局状态                 \\
学习曲线       & 较陡（函数式思维）        & 较平缓（命令式直觉）          \\
\bottomrule
\end{tabular}
\end{table}

% =============================================
\section{完整示例：用 JAX 训练线性回归}
% =============================================

下面是一个完整的线性回归训练示例，展示了 JAX 四大变换的综合运用：

\begin{lstlisting}
import jax
import jax.numpy as jnp

# 1. 生成合成数据
key = jax.random.PRNGKey(0)
key, subkey = jax.random.split(key)
X = jax.random.normal(subkey, (100, 3))
w_true = jnp.array([1.0, 2.0, -3.0])
y = X @ w_true + 0.1 * jax.random.normal(
    jax.random.split(key)[1], (100,))

# 2. 定义模型与损失函数
def predict(params, x):
    return jnp.dot(x, params)

def loss(params, x, y):
    pred = predict(params, x)
    return jnp.mean((pred - y) ** 2)

# 3. 定义训练步骤（jit 编译加速）
@jax.jit
def train_step(params, x, y, lr=0.01):
    grads = jax.grad(loss)(params, x, y)
    return params - lr * grads

# 4. 训练循环
params = jnp.zeros(3)
for epoch in range(100):
    params = train_step(params, X, y)
    if epoch % 20 == 0:
        print(f"Epoch {epoch}: loss = {loss(params, X, y):.4f}")

print(f"Learned weights: {params}")
# 期望输出接近 [1.0, 2.0, -3.0]
\end{lstlisting}

% =============================================
\section{学习资源}
% =============================================

\begin{itemize}
    \item \textbf{官方文档}：\url{https://jax.readthedocs.io}
    \item \textbf{JAX 101 教程}：官方提供的从零开始的系列教程
    \item \textbf{JAX in Action}： Manning 出版社的实战书籍
    \item \textbf{DeepMind 的 JAX 生态}：\\ \url{https://www.deepmind.com/blog/using-jax-to-accelerate-our-research}
\end{itemize}

% =============================================
\section{总结}
% =============================================

JAX 以\textbf{函数变换}为核心抽象，将自动微分、编译加速和并行计算统一在一个简洁而强大的框架中。虽然其纯函数式设计范式需要一定的学习成本，但一旦掌握，便能以极高的表达力和组合性来构建复杂的机器学习研究流水线。

对于追求计算效率、需要高阶微分或希望探索新型模型架构的研究者而言，JAX 是一个值得深入学习的工具。

\end{document}

这是一篇完整的 JAX 框架入门介绍，主要内容包括：

1. **什么是 JAX**：核心定位与优势概述
2. **安装与快速上手**：最小可运行示例
3. **四大核心变换**：`grad`、`jit`、`vmap`、`pmap` 的详细介绍与代码示例
4. **纯函数式约束**：不可变数组、显式随机数管理等常见"坑"
5. **生态系统**：Flax、Optax 等上层库的介绍
6. **与 PyTorch 对比**：帮助有 PyTorch 背景的读者快速理解差异
7. **完整示例**：线性回归训练的端到端代码

你可以直接用 `xelatex` 编译生成 PDF。如需调整章节深度、增减内容或修改排版样式，随时告诉我。